\documentclass[UTF8,11pt,a4paper]{ctexart}
\newcommand{\DocTitle}{决策策略说明}
\newcommand{\DocSubtitle}{从开局采购到回放复盘}
% Shared layout for the CW2 companion documents.
\usepackage[margin=2.1cm,headheight=16pt]{geometry}
\usepackage{amsmath,amssymb,booktabs,array,tabularx,ragged2e}
\usepackage{enumitem,xcolor,fancyhdr,listings,tcolorbox}
\tcbuselibrary{breakable}
\usepackage{hyperref}
\usepackage{bookmark}

\newcommand{\DocVersion}{3.0}
\newcommand{\EngineVersion}{1.1}
\newcommand{\DocDate}{2026-11-23}
\definecolor{ink}{HTML}{1F2937}
\definecolor{accent}{HTML}{92400E}
\definecolor{muted}{HTML}{4B5563}
\definecolor{linkink}{HTML}{174678}
\definecolor{codebg}{HTML}{F5F7FA}
\hypersetup{unicode,colorlinks=true,linkcolor=linkink,urlcolor=linkink,
  filecolor=linkink,pdftitle={SOE104 CW2 \DocTitle\space v\DocVersion},
  pdfauthor={Group A--E}}
\ctexset{section={format={\Large\bfseries\color{accent}}},
  subsection={format={\large\bfseries\color{ink}}}}
\pagestyle{fancy}
\fancyhf{}
\fancyhead[L]{\small\color{muted}SOE104 CW2 · Oregon Trail}
\fancyhead[R]{\small\color{muted}\DocTitle\ v\DocVersion}
\fancyfoot[C]{\small\color{muted}\thepage}
\renewcommand{\headrulewidth}{0.4pt}
\setlength{\parskip}{0.35em}
\setlength{\emergencystretch}{1.5em}
\renewcommand{\arraystretch}{1.23}
\setlist{itemsep=3pt,topsep=3pt,parsep=0pt,leftmargin=1.7em}
\newcolumntype{Y}{>{\RaggedRight\arraybackslash}X}
\newcolumntype{L}[1]{>{\RaggedRight\arraybackslash}p{#1}}
\newtcolorbox{keybox}[1]{breakable,title={#1},colback=orange!4,
  colframe=accent,boxrule=0.55pt,arc=2pt,fonttitle=\bfseries}
\newtcolorbox{warnbox}[1]{breakable,title={#1},colback=red!3,
  colframe=red!45!black,boxrule=0.55pt,arc=2pt,fonttitle=\bfseries}
\lstset{basicstyle=\small\ttfamily,backgroundcolor=\color{codebg},
  xleftmargin=0.9em,frame=l,rulecolor=\color{linkink!45},
  breaklines=true,columns=fullflexible,keepspaces=true,showstringspaces=false,
  aboveskip=0.65em,belowskip=0.65em}
\newcommand{\code}[1]{\texttt{#1}}
\newcommand{\source}[1]{\textup{\small\textcolor{muted}{[#1]}}}
\newcommand{\docref}[2]{\href{./#1.pdf}{#2}}
\renewcommand{\maketitle}{%
  \begin{center}
    {\LARGE\bfseries CW2 \DocTitle\par}
    \vspace{0.35em}{\large\DocSubtitle\par}
    \vspace{0.75em}{\small\color{muted}Group A--E · 整理：A
      \quad\DocDate\quad 文档 v\DocVersion\par}
    \vspace{0.2em}{\small\color{muted}Engine / Marking v\EngineVersion\par}
  \end{center}
  \vspace{0.4em}}


\begin{document}
\maketitle
\thispagestyle{fancy}

\begin{keybox}{这套策略想做到什么}
120 回合内走到 2000 英里，尽量让队员健康到达，车上还有物资。
出发前准备一套能运转的队伍和补给。上路后，每回合根据实际的血量、存粮、牛数和路程重新判断。
\end{keybox}

\section{从出发到终点}

\begin{center}\small
出发前配队和采购
\quad$\longrightarrow$\quad
每回合选配速和口粮、处理事件、分配动作
\quad$\longrightarrow$\quad
终局计分
\end{center}

\subsection{出发前准备一个能持续的组合}

先用猎人、医生、机械师和侦察员组成一支四人队。猎人补粮，医生降低事件伤害，
机械师负责修车和防守，侦察员补里程。多带一名队员每回合多一个动作，也多一份口粮。
四人队刚好能把这些工作分开，便于看出哪个环节需要改。

开局买够基本口粮、运力和弹药。四头牛已经能让车队按完整配速行进。
第五头牛可以用来缓冲一次损失，但也要多吃粮。多买弹药可以增加途中狩猎的机会，
前提是猎人有空打猎，车里也装得下。两种采购思路都可以用回放对照。

\subsection{上路后用同一套判断顺序}

\begin{enumerate}
  \item \textbf{先检查能不能承受后果。}买不起的方案先排除。
  再看失败时有没有人会死、下一餐够不够。所有方案都有危险时，优先降低全队覆灭的风险。
  \item \textbf{看看还赶不赶得上。}把剩下的路和回合放在一起看，
  算上牛拉车和队员探路，也留一点绕道余量。牛少了，或者有人本来要探路、现在必须打猎，估计就跟着改。
  \item \textbf{在可行方案里挑代价小的。}比较会丢多少里程、血量和物资，也看冒险失败的后果。
  里程充裕时，可以用一点时间保住血量或补给。接近截止回合时，额外的几英里可能更重要。
\end{enumerate}

比如掉血之后，治疗会耗粮，还占用一次原本可以拿来赶路的行动。
因此比较绕路和受伤时，我们会把后续恢复的代价也考虑进去。具体怎么换算留到完整版。

\clearpage
\section{跟着一回合走一遍}

\noindent 假设现在四个人、四头牛，车里有 34 单位食物。
队里有人受了伤，但还能行动。猎人有弹药，行程也留有余量。下面是一种可能的选择。

\subsection{回合开始：同时选配速和口粮}

选正常配速和普通口粮（\code{MEAGER}）。这时没有赶路压力，提速会多掉血，还有失牛的风险。
现有粮食能覆盖普通档这一餐，猎人也还能补给。粮紧时可以考虑更低的口粮，
同时安排补给，因为长期压低口粮会让队伍掉血。

\subsection{行进和喂食：先由引擎结算}

队员先吃，四个人用掉 12 单位食物。四头牛接着吃掉 8 单位。车里剩下 14 单位。
我们既关心队员吃没吃饱，也关心全车一共消耗了多少。前者影响当回合饥饿，
后者决定后面还能撑几轮。

\subsection{遇到事件：把后果放回行程里比较}

假设遇到河流。绕行会损失 5 英里，剩余行程还有余量，可以绕行把现金留给采购。
如果同样的 5 英里会让队伍赶不上终点，而且付得起船费，付费过河就有更高价值。
这里的差别来自队伍当时的进度。遇到会扣粮的事件，也用同样的办法看后果。
本轮已经吃过了，少掉的食物要算进下一餐和接下来能做的动作。

\subsection{事件之后：用成员动作补上缺口}

在这个例子里，猎人有弹药和仓位，可以打猎带回 10 单位食物，车里变成 24。
接着，医生如果判断值得治疗，可以花 3 单位，仍剩 21。侦察员探路补回推进里程。
机械师根据是否受威胁和队员伤势决定防守还是帮忙赶路。

\begin{keybox}{为什么每一轮都要重新看}
这轮结束时有 21 单位食物。如果下轮仍是四人四牛、普通口粮，够付全车 20 单位。
如果猎人这一轮没能打猎，即便医生省下治疗用的粮，也只有 14 单位。
这够四名队员吃普通口粮，却不够全车，下一轮要重新考虑补给。
同一个配速和口粮组合，适不适合要看最新状态。
\end{keybox}

\clearpage
\section{这些决定如何接起来}

项目把每回合的决定放在三个时点。回合开始的策略能看全队的状态，负责定配速和口粮。
事件处理拿到的是事件本身，车况要参考策略留下的本轮记录。
行进和喂食已经发生，所以这份记录里要区分\textbf{开头看到的值}和
\textbf{推算到事件时的值}。成员在事件结算后再行动，用自己收到的当时信息来决定。

\subsection{成员有默认职责，也可以换动作}

\begin{center}\small
\begin{tabularx}{\linewidth}{@{}L{1.8cm}Y Y@{}}
\toprule
角色 & 平时主要做什么 & 什么时候调整 \\
\midrule
猎人 & 补足后面要吃的粮。 & 粮食快满时，再打猎可能只进仓一点，弹药照扣。可改为推进。 \\
医生 & 活着就能减轻全队受到的事件伤害，治疗还能照顾最伤的队员。 &
本轮吃饱、粮食够用且治疗有收益时才安排治疗。先留够下一餐。 \\
机械师 & 存活时，安全修理故障可以免耗零件。遇险时由他判断是否防守。 &
整队防守每回合只生效一次。没有防守需要时，可以帮忙推进。 \\
侦察员 & 探路多走些里程，给后续绕道留空间。 &
里程已有余量而补给或保命更紧迫时，可以改做其他动作。 \\
\bottomrule
\end{tabularx}
\end{center}

\subsection{越靠近终点，越看重能留下什么}

路还长时，食物和弹药主要用来支撑下一段路，现金用来交易或付费过关。
如果已经确认能赶到终点，就多留意最终血量和存货。
终局只把食物、牛、弹药和零件计作货物，现金留在钱包里没有货物分。

\subsection{用回放看设计有没有解决问题}

用同一组平台种子对照改动前后。先看有没有到达，再看哪一回合开始缺粮、谁受伤或死亡、
牛有没有掉，以及最终留下多少货物。比如想验证猎人是否应该更早补粮，
先只改狩猎时机，核对饥饿回合是否减少、额外打猎是否影响里程，再决定是否保留改动。

\medskip
\noindent\small 这是团队拟采用并用回放校准的设计。现有 Java 仍是基础版。
事件数值、推导和代码接口见
\docref{04_decision_model}{《04 决策模型》完整版}。

\end{document}
